\documentclass[12pt,a4paper]{report}

% ============================================
% PACKAGES
% ============================================

% Font Roboto
\usepackage[sfdefault]{roboto}
\usepackage[T1]{fontenc}
\usepackage[utf8]{inputenc}

% Lingua italiana
\usepackage[italian]{babel}

% Geometria pagina
\usepackage[margin=2.5cm]{geometry}

% Colori e grafica
\usepackage{xcolor}
\usepackage{graphicx}
\usepackage{tikz}
\usetikzlibrary{shapes.geometric, arrows.meta, positioning, fit, backgrounds, calc}

% Codice sorgente
\usepackage{listings}
\usepackage{inconsolata}

% Tabelle avanzate
\usepackage{booktabs}
\usepackage{tabularx}
\usepackage{longtable}

% Link e PDF metadata
\usepackage{hyperref}
\hypersetup{
    colorlinks=true,
    linkcolor=blue!60!black,
    urlcolor=blue!70!black,
    citecolor=green!50!black,
    pdftitle={Open Data Chunker - Documentazione Tecnica},
    pdfauthor={Gabriele},
    pdfsubject={ETL Pipeline per RNA Open Data}
}

% Caption
\usepackage{caption}
\captionsetup{font=small, labelfont=bf}

% Quote e box
\usepackage{tcolorbox}
\tcbuselibrary{skins, breakable}

% ============================================
% DEFINIZIONI COLORI
% ============================================
\definecolor{codebackground}{RGB}{248, 248, 248}
\definecolor{codekeyword}{RGB}{0, 119, 170}
\definecolor{codestring}{RGB}{162, 35, 39}
\definecolor{codecomment}{RGB}{87, 166, 74}
\definecolor{primaryblue}{RGB}{41, 98, 255}
\definecolor{accentgreen}{RGB}{0, 200, 83}
\definecolor{warnorange}{RGB}{255, 152, 0}

% ============================================
% STILE LISTINGS
% ============================================
\lstdefinestyle{codestyle}{
    backgroundcolor=\color{codebackground},
    basicstyle=\ttfamily\small,
    keywordstyle=\color{codekeyword}\bfseries,
    stringstyle=\color{codestring},
    commentstyle=\color{codecomment}\itshape,
    breakatwhitespace=false,
    breaklines=true,
    captionpos=b,
    keepspaces=true,
    numbers=left,
    numbersep=8pt,
    numberstyle=\tiny\color{gray},
    showspaces=false,
    showstringspaces=false,
    showtabs=false,
    tabsize=4,
    frame=single,
    framerule=0.5pt,
    rulecolor=\color{gray!30}
}
\lstset{style=codestyle}

\lstdefinelanguage{docker}{
    keywords={FROM, RUN, COPY, WORKDIR, ENV, CMD, EXPOSE, VOLUME, ARG, ENTRYPOINT},
    keywordstyle=\color{codekeyword}\bfseries,
    sensitive=true,
    comment=[l]{\#},
    morestring=[b]",
    morestring=[b]'
}

% ============================================
% TCOLORBOX STYLES
% ============================================
\newtcolorbox{infobox}[1][]{
    colback=primaryblue!5,
    colframe=primaryblue,
    fonttitle=\bfseries,
    title=#1,
    breakable
}

\newtcolorbox{tipbox}[1][]{
    colback=accentgreen!5,
    colframe=accentgreen!80!black,
    fonttitle=\bfseries,
    title=#1,
    breakable
}

% ============================================
% TIKZ STYLES
% ============================================
\tikzstyle{component} = [rectangle, rounded corners, minimum width=3cm, minimum height=1cm, text centered, draw=black, fill=primaryblue!20]
\tikzstyle{arrow} = [thick, ->, >=Stealth]

% Database icon command
\newcommand{\dbicon}[4]{%
    % #1=x, #2=y, #3=label, #4=fill color
    \begin{scope}[shift={(#1,#2)}]
        \fill[#4, draw=black, rounded corners=1pt] (-0.8,-0.4) rectangle (0.8,0.4);
        \draw[black] (-0.8,0) -- (0.8,0);
        \draw[black] (-0.8,-0.2) -- (0.8,-0.2);
        \fill[#4, draw=black] (-0.8,0.4) arc (180:360:0.8 and 0.2) -- (0.8,0.4) arc (0:180:0.8 and 0.2) -- cycle;
        \draw[black] (-0.8,0.4) arc (180:360:0.8 and 0.2);
        \node at (0,-0.6) {\small #3};
    \end{scope}
}

% ============================================
% DOCUMENTO
% ============================================

\begin{document}

% Frontespizio
\begin{titlepage}
    \centering
    \vspace*{2cm}
    
    {\Huge\bfseries Open Data Chunker\par}
    \vspace{0.5cm}
    {\Large\itshape Pipeline ETL ad Alte Prestazioni\par}
    
    \vspace{2cm}
    
    \begin{tikzpicture}[scale=0.8, transform shape]
        \node[component, fill=warnorange!30] (xml) at (0,0) {XML Files};
        \node[component] (parser) at (4,0) {Parser};
        \dbicon{8}{0}{Parquet}{accentgreen!30}
        \node[component, fill=primaryblue!30] (query) at (12,0) {Query Engine};
        
        \draw[arrow] (xml) -- (parser);
        \draw[arrow] (parser) -- (7.2,0);
        \draw[arrow] (8.8,0) -- (query);
    \end{tikzpicture}
    
    \vspace{2cm}
    
    {\large Documentazione Tecnica\par}
    \vspace{1cm}
    {\large Versione 1.0\par}
    
    \vfill
    
    {\large \today\par}
\end{titlepage}

% Indice
\tableofcontents
\newpage

% Capitoli
\chapter{Introduzione}

\section{Contesto}

Il \textbf{Registro Nazionale degli Aiuti} (RNA) è la banca dati nazionale istituita presso il Ministero dello Sviluppo Economico che raccoglie le informazioni relative agli aiuti di Stato e agli aiuti \textit{de minimis} concessi in Italia. L'RNA pubblica periodicamente dataset in formato \textbf{Open Data XML}, consentendo l'accesso pubblico a informazioni dettagliate su milioni di contributi erogati.

Questi dataset presentano caratteristiche che li rendono complessi da gestire:

\begin{itemize}
    \item \textbf{Volumi elevati}: centinaia di file XML, ciascuno contenente decine di migliaia di record
    \item \textbf{Struttura gerarchica}: dati annidati su più livelli (Aiuti $\rightarrow$ Componenti $\rightarrow$ Strumenti)
    \item \textbf{Qualità variabile}: presenza di caratteri XML non validi e encoding inconsistente
    \item \textbf{Assenza di formati analitici}: i dati sono distribuiti solo in XML, non ottimale per query analitiche
\end{itemize}

\section{Obiettivi del Progetto}

\textbf{Open Data Chunker} è una pipeline ETL (\textit{Extract, Transform, Load}) progettata per rispondere a queste sfide. Gli obiettivi principali sono:

\begin{enumerate}
    \item \textbf{Parsing robusto}: capacità di processare file XML corrotti o malformati senza interruzioni
    \item \textbf{Performance scalabili}: elaborazione parallela multi-processo per sfruttare CPU multi-core
    \item \textbf{Output analitico}: conversione in formato \textbf{Apache Parquet} con partizionamento temporale
    \item \textbf{Query interattive}: integrazione con motori SQL (\textbf{DuckDB}) per interrogazioni ad-hoc
    \item \textbf{Riproducibilità}: containerizzazione Docker per ambienti consistenti
\end{enumerate}

\begin{infobox}[Punti di Forza]
\begin{itemize}
    \item Pipeline end-to-end: dal raw XML al dataset queryable
    \item Fault-tolerant: gestione automatica degli errori
    \item Compressione 10x rispetto all'XML originale
    \item Deploy immediato via Docker
\end{itemize}
\end{infobox}

\section{Target Audience}

Questa documentazione è rivolta a \textbf{Software Engineer} e \textbf{Data Engineer} interessati a:

\begin{itemize}
    \item Comprendere l'architettura di una pipeline ETL Python moderna
    \item Valutare pattern di parallelismo e gestione memoria
    \item Replicare o estendere la soluzione per casi d'uso simili
\end{itemize}

\chapter{Architettura}

\section{Panoramica}

L'architettura di Open Data Chunker segue un pattern \textbf{layered} con separazione netta delle responsabilità. Il sistema è composto da quattro layer principali, illustrati nel diagramma seguente.

\begin{figure}[h]
\centering
\begin{tikzpicture}[node distance=1.5cm, scale=0.9, transform shape]
    
    % Stile database migliorato
    \tikzstyle{dbbox} = [rectangle, rounded corners=3pt, minimum width=3.5cm, minimum height=1.2cm, text centered, draw=black, fill=accentgreen!25, font=\small]
    
    % Input Layer
    \node[component, fill=warnorange!30, minimum width=4cm] (xml) {Input Layer\\{\small XML Files}};
    
    % Processing Layer
    \node[component, fill=primaryblue!20, minimum width=4cm, right=2cm of xml] (process) {Processing Layer\\{\small lxml + Sanitizer}};
    
    % Storage Layer - Box stile database
    \node[dbbox, right=2cm of process] (storage) {Storage Layer\\{\small Apache Parquet}};
    \draw[black, thick] ([yshift=-3pt]storage.north west) -- ([yshift=-3pt]storage.north east);
    \draw[black, thick] ([yshift=-8pt]storage.north west) -- ([yshift=-8pt]storage.north east);
    
    % Query Layer
    \node[component, fill=primaryblue!40, minimum width=4cm, below=2cm of storage] (query) {Query Layer\\{\small DuckDB + Polars}};
    
    % CLI
    \node[component, fill=gray!20, minimum width=3cm, below=2cm of process] (cli) {CLI\\{\small Click}};
    
    % Arrows
    \draw[arrow] (xml) -- node[above, font=\small] {stream} (process);
    \draw[arrow] (process) -- node[above, font=\small] {batch write} (storage);
    \draw[arrow] (storage) -- (query);
    \draw[arrow] (cli) -- (process);
    \draw[arrow] (cli) -- (query);
    
\end{tikzpicture}
\caption{Architettura a layer di Open Data Chunker}
\label{fig:architecture}
\end{figure}

\section{Layer di Input}

Il layer di input gestisce i file XML grezzi provenienti dall'RNA Open Data. Caratteristiche:

\begin{itemize}
    \item Supporto per file singoli o directory ricorsive
    \item Discovery automatico tramite glob pattern \texttt{*.xml}
    \item Nessun caricamento in memoria dell'intero file (streaming)
\end{itemize}

\section{Layer di Processing}

Il core della pipeline, implementato in \texttt{src/parser.py}. Composto da:

\begin{table}[h]
\centering
\begin{tabular}{@{}lp{8cm}@{}}
\toprule
\textbf{Componente} & \textbf{Responsabilità} \\
\midrule
\texttt{CleanFileInputStream} & Wrapper che filtra caratteri XML non validi in tempo reale \\
\texttt{iterparse()} & Parser SAX-like per elaborazione streaming \\
\texttt{ProcessPoolExecutor} & Orchestratore del parallelismo multi-processo \\
\texttt{flush\_batches()} & Writer batch su Parquet con partizionamento \\
\bottomrule
\end{tabular}
\caption{Componenti del Processing Layer}
\end{table}

\section{Layer di Storage}

Output in formato \textbf{Apache Parquet} con le seguenti caratteristiche:

\begin{itemize}
    \item \textbf{Columnar storage}: compressione ottimale e predicate pushdown
    \item \textbf{Hive partitioning}: directory \texttt{ANNO=YYYY} per partition pruning
    \item \textbf{Schema tipizzato}: definizioni PyArrow in \texttt{src/models.py}
\end{itemize}

\begin{tipbox}[Vantaggio: Partition Pruning]
Le query filtrate per anno leggono solo le partizioni rilevanti, riducendo I/O fino al 90\% su dataset multi-anno.
\end{tipbox}

\section{Layer di Query}

Dual-engine per flessibilità:

\begin{description}
    \item[DuckDB] Motore SQL analitico in-process. Ideale per aggregazioni e join complessi via CLI.
    \item[Polars] DataFrame library lazy-first. Utilizzato per export streaming memory-efficient.
\end{description}

\section{Flusso dei Dati}

\begin{enumerate}
    \item L'utente invoca \texttt{parse --input data/} via CLI
    \item Il CLI discovery tutti i file \texttt{*.xml} nella directory
    \item I file vengono distribuiti tra N worker processes
    \item Ogni worker:
    \begin{enumerate}
        \item Apre il file tramite \texttt{CleanFileInputStream}
        \item Itera gli elementi \texttt{<AIUTO>} con \texttt{iterparse()}
        \item Estrae i campi in batch da 10.000 record
        \item Scrive il batch su Parquet partizionato
    \end{enumerate}
    \item Al termine, statistiche aggregate vengono loggite
\end{enumerate}

\chapter{Implementazione}

\section{Stack Tecnologico}

La soluzione utilizza un moderno stack Python per data engineering:

\begin{table}[h]
\centering
\begin{tabular}{@{}llp{6cm}@{}}
\toprule
\textbf{Libreria} & \textbf{Versione} & \textbf{Utilizzo} \\
\midrule
\texttt{lxml} & 5.x & Parsing XML streaming \\
\texttt{polars} & 1.x & DataFrame processing \\
\texttt{pyarrow} & 18.x & Schema definition e Parquet I/O \\
\texttt{duckdb} & 1.x & Query engine SQL \\
\texttt{click} & 8.x & CLI framework \\
\texttt{tqdm} & 4.x & Progress bar \\
\bottomrule
\end{tabular}
\caption{Dipendenze principali}
\end{table}

\section{Sanitizzazione XML}

Un problema critico dei dataset RNA è la presenza di \textbf{caratteri XML non validi}. Questi includono:

\begin{itemize}
    \item Caratteri di controllo ASCII: \texttt{0x00-0x08}, \texttt{0x0B-0x0C}, \texttt{0x0E-0x1F}
    \item Entità numeriche malformate: \texttt{\&\#0;} - \texttt{\&\#31;}
\end{itemize}

La classe \texttt{CleanFileInputStream} risolve questo problema tramite un wrapper file-like che filtra lo stream in tempo reale:

\begin{lstlisting}[language=Python, caption=Sanitizzazione streaming]
class CleanFileInputStream:
    def __init__(self, filename):
        self.f = open(filename, 'rb')
        self.invalid_xml_chars = re.compile(
            b'[\x00-\x08\x0b\x0c\x0e-\x1f]|'
            b'&#(?:[0-8]|1[1-2]|1[4-9]|2[0-9]|3[0-1]);'
        )

    def read(self, size=-1):
        data = self.f.read(size)
        return self.invalid_xml_chars.sub(b'', data)
\end{lstlisting}

\begin{tipbox}[Design Choice]
Il pattern wrapper consente di integrare la sanitizzazione senza modificare il codice di parsing. Il parser \texttt{lxml} riceve uno stream già pulito.
\end{tipbox}

\section{Parsing Streaming}

Per gestire file multi-gigabyte senza esaurire la memoria, utilizziamo \texttt{lxml.etree.iterparse()} con cleanup aggressivo:

\begin{lstlisting}[language=Python, caption=Iterparse con memory management]
context = etree.iterparse(
    clean_stream, 
    events=("end",), 
    tag=f"{NS}AIUTO"
)

for event, elem in context:
    # Estrazione dati...
    
    # Release memory
    elem.clear()
    while elem.getprevious() is not None:
        del elem.getparent()[0]
\end{lstlisting}

Il pattern \texttt{elem.clear()} + deletion dei siblings previene memory leak tipici di \texttt{iterparse}.

\section{Parallelismo Multi-Processo}

L'elaborazione parallela sfrutta \texttt{ProcessPoolExecutor} per bypassare il GIL:

\begin{lstlisting}[language=Python, caption=Orchestrazione parallela]
with ProcessPoolExecutor(max_workers=workers) as executor:
    futures = {
        executor.submit(process_file, f, output): f 
        for f in files
    }
    
    for future in as_completed(futures):
        stats = future.result()
        # Aggregate statistics...
\end{lstlisting}

\begin{figure}[h]
\centering
\begin{tikzpicture}[node distance=0.8cm, scale=0.85, transform shape]
    % Stile file parquet
    \tikzstyle{parquetfile} = [rectangle, rounded corners=2pt, minimum width=2.2cm, minimum height=0.8cm, text centered, draw=black, fill=accentgreen!25, font=\scriptsize]
    
    \node[component, minimum width=3cm] (main) {Main Process};
    
    \node[component, fill=accentgreen!20, below left=1.5cm and 1cm of main] (w1) {Worker 1};
    \node[component, fill=accentgreen!20, below=1.5cm of main] (w2) {Worker 2};
    \node[component, fill=accentgreen!20, below right=1.5cm and 1cm of main] (w3) {Worker N};
    
    \node[above right=0.3cm and -0.5cm of w2] {...};
    
    \draw[arrow] (main) -- (w1);
    \draw[arrow] (main) -- (w2);
    \draw[arrow] (main) -- (w3);
    
    \node[parquetfile, below=1cm of w1] (p1) {file1.parquet};
    \node[parquetfile, below=1cm of w2] (p2) {file2.parquet};
    \node[parquetfile, below=1cm of w3] (p3) {fileN.parquet};
    
    \draw[arrow] (w1) -- (p1);
    \draw[arrow] (w2) -- (p2);
    \draw[arrow] (w3) -- (p3);
\end{tikzpicture}
\caption{Modello fork-join per elaborazione parallela}
\end{figure}

\section{Batching e Scrittura}

Per ottimizzare le scritture su disco, i record vengono accumulati in batch di 10.000 elementi prima del flush:

\begin{lstlisting}[language=Python, caption=Flush batch su Parquet]
BATCH_SIZE = 10000

if len(batch_aiuti) >= BATCH_SIZE:
    pq.write_to_dataset(
        table,
        root_path=str(base_path / "aiuti"),
        partition_cols=['ANNO'],
        existing_data_behavior='overwrite_or_ignore'
    )
\end{lstlisting}

Il parametro \texttt{partition\_cols=['ANNO']} abilita il \textbf{Hive-style partitioning}, creando automaticamente directory \texttt{ANNO=2022/}, \texttt{ANNO=2023/}, etc.

\chapter{Modello Dati}

\section{Struttura Gerarchica XML}

I file XML dell'RNA presentano una struttura gerarchica a tre livelli:

\begin{figure}[h]
\centering
\begin{tikzpicture}[
    level 1/.style={sibling distance=5cm, level distance=1.5cm},
    level 2/.style={sibling distance=3cm, level distance=1.5cm},
    every node/.style={rectangle, rounded corners, draw, fill=primaryblue!10, minimum width=2.5cm, minimum height=0.8cm, font=\small}
]
\node {AIUTO}
    child {node {COMPONENTE\_1}
        child {node {STRUM\_1}}
        child {node {STRUM\_2}}
    }
    child {node {COMPONENTE\_2}
        child {node {STRUM\_3}}
    };
\end{tikzpicture}
\caption{Struttura gerarchica dei dati RNA}
\end{figure}

\section{Normalizzazione in Tre Tabelle}

Il parser trasforma la struttura gerarchica in tre tabelle relazionali normalizzate:

\subsection{Tabella AIUTI}

Entità principale contenente le informazioni sul contributo:

\begin{table}[h]
\centering
\small
\begin{tabular}{@{}lll@{}}
\toprule
\textbf{Campo} & \textbf{Tipo} & \textbf{Descrizione} \\
\midrule
CAR & string & Codice Aiuto Regionale (PK logica) \\
COR & string & Codice Operazione Registro \\
TITOLO\_MISURA & string & Denominazione del bando \\
DATA\_CONCESSIONE & string & Data di erogazione \\
DENOMINAZIONE\_BENEFICIARIO & string & Ragione sociale beneficiario \\
CODICE\_FISCALE\_BENEFICIARIO & string & CF/P.IVA \\
REGIONE\_BENEFICIARIO & string & Regione di appartenenza \\
CUP & string & Codice Unico Progetto \\
ANNO & int32 & Partizione temporale \\
FILE\_SOURCE & string & Tracciabilità file sorgente \\
\bottomrule
\end{tabular}
\caption{Schema tabella AIUTI}
\end{table}

\subsection{Tabella COMPONENTI}

Dettaglio delle componenti di aiuto, in relazione 1:N con AIUTI:

\begin{table}[h]
\centering
\small
\begin{tabular}{@{}lll@{}}
\toprule
\textbf{Campo} & \textbf{Tipo} & \textbf{Descrizione} \\
\midrule
ID\_COMPONENTE\_AIUTO & string & Identificativo componente (PK) \\
CAR\_AIUTO & string & FK verso AIUTI.CAR \\
COR\_AIUTO & string & FK verso AIUTI.COR \\
COD\_REGOLAMENTO & string & Regolamento UE applicato \\
SETTORE\_ATTIVITA & string & Codice ATECO \\
ANNO & int32 & Partizione \\
\bottomrule
\end{tabular}
\caption{Schema tabella COMPONENTI}
\end{table}

\subsection{Tabella STRUMENTI}

Strumenti finanziari applicati, in relazione 1:N con COMPONENTI:

\begin{table}[h]
\centering
\small
\begin{tabular}{@{}lll@{}}
\toprule
\textbf{Campo} & \textbf{Tipo} & \textbf{Descrizione} \\
\midrule
ID\_COMPONENTE\_AIUTO & string & FK verso COMPONENTI \\
COD\_STRUMENTO & string & Codice strumento finanziario \\
IMPORTO\_NOMINALE & float64 & Valore nominale contributo \\
ELEMENTO\_DI\_AIUTO & float64 & Equivalente sovvenzione \\
ANNO & int32 & Partizione \\
\bottomrule
\end{tabular}
\caption{Schema tabella STRUMENTI}
\end{table}

\section{Diagramma ER}

\begin{figure}[h]
\centering
\begin{tikzpicture}[
    entity/.style={rectangle, draw, fill=primaryblue!20, minimum width=3cm, minimum height=1.5cm, font=\bfseries},
    relationship/.style={diamond, draw, fill=warnorange!30, minimum width=1.5cm, aspect=2},
    scale=0.9, transform shape
]
    \node[entity] (aiuti) at (0,0) {AIUTI};
    \node[entity] (comp) at (5,0) {COMPONENTI};
    \node[entity] (strum) at (10,0) {STRUMENTI};
    
    \node[relationship] (r1) at (2.5,0) {1:N};
    \node[relationship] (r2) at (7.5,0) {1:N};
    
    \draw[-] (aiuti) -- (r1);
    \draw[-] (r1) -- (comp);
    \draw[-] (comp) -- (r2);
    \draw[-] (r2) -- (strum);
    
    \node[below=0.3cm of aiuti, font=\small\ttfamily] {CAR, COR};
    \node[below=0.3cm of comp, font=\small\ttfamily] {ID\_COMPONENTE};
    \node[below=0.3cm of strum, font=\small\ttfamily] {IMPORTO, ELEMENTO};
\end{tikzpicture}
\caption{Diagramma ER delle tabelle normalizzate}
\end{figure}

\section{Partizionamento Hive-Style}

L'output è organizzato con \textbf{Hive-style partitioning} sulla colonna \texttt{ANNO}:

\begin{lstlisting}[language=bash, caption=Struttura directory output]
public/parquet/
    aiuti/
        ANNO=2022/
            part-0.parquet
            part-1.parquet
        ANNO=2023/
            part-0.parquet
    componenti/
        ANNO=2022/
        ANNO=2023/
    strumenti/
        ANNO=2022/
        ANNO=2023/
\end{lstlisting}

Questo layout consente:
\begin{itemize}
    \item \textbf{Partition pruning}: query filtrate per anno leggono solo le partizioni rilevanti
    \item \textbf{Parallelismo}: file distinti possono essere letti concorrentemente
    \item \textbf{Gestibilità}: facile eliminazione o aggiornamento di singoli anni
\end{itemize}

\chapter{Performance}

\section{Metriche di Riferimento}

Le seguenti metriche sono state misurate su un dataset reale dell'RNA composto da circa 165 file XML:

\begin{table}[h]
\centering
\begin{tabular}{@{}ll@{}}
\toprule
\textbf{Metrica} & \textbf{Valore} \\
\midrule
Velocità di parsing & $\sim$50.000 record/sec (8 worker) \\
Utilizzo memoria per worker & $\sim$200 MB \\
Rapporto compressione & $\sim$10x vs XML grezzo \\
\bottomrule
\end{tabular}
\caption{Benchmark su dataset RNA reale}
\end{table}

\section{Scalabilità Orizzontale}

La performance scala linearmente con il numero di worker fino al numero di core fisici disponibili:

\begin{figure}[h]
\centering
\begin{tikzpicture}
    \begin{axis}[
        width=10cm,
        height=6cm,
        xlabel={Numero Worker},
        ylabel={Throughput (record/sec)},
        xmin=1, xmax=8,
        ymin=0, ymax=55000,
        xtick={1,2,4,6,8},
        ytick={0,10000,20000,30000,40000,50000},
        legend pos=north west,
        grid=major,
        grid style={gray!30}
    ]
    \addplot[color=primaryblue, mark=*, thick] coordinates {
        (1, 7000)
        (2, 14000)
        (4, 27000)
        (6, 40000)
        (8, 50000)
    };
    \addlegendentry{Misurato}
    
    \addplot[color=gray, dashed] coordinates {
        (1, 7000)
        (8, 56000)
    };
    \addlegendentry{Lineare ideale}
    \end{axis}
\end{tikzpicture}
\caption{Scalabilità throughput vs numero worker}
\end{figure}

\begin{infobox}[Nota sulle Performance]
Il leggero scostamento dalla linearità ideale è dovuto a:
\begin{itemize}
    \item Overhead di serializzazione/deserializzazione tra processi
    \item Contesa I/O su disco per scritture Parquet concorrenti
    \item Scheduling overhead del sistema operativo
\end{itemize}
\end{infobox}

\section{Compressione Parquet}

Il formato Parquet offre vantaggi significativi rispetto all'XML:

\begin{table}[h]
\centering
\begin{tabular}{@{}lll@{}}
\toprule
\textbf{Aspetto} & \textbf{XML} & \textbf{Parquet} \\
\midrule
Storage size & 1 GB & $\sim$100 MB \\
Tempo lettura completa & $>$60s & $<$5s \\
Query selettive & Full scan & Partition pruning \\
Tipizzazione & Assente & Nativa \\
\bottomrule
\end{tabular}
\caption{Confronto XML vs Parquet}
\end{table}

\section{Ottimizzazioni Chiave}

\subsection{Streaming Parser}

L'utilizzo di \texttt{iterparse()} evita di caricare l'intero DOM in memoria:

\begin{itemize}
    \item Memory footprint costante indipendentemente dalla dimensione file
    \item Possibilità di processare file multi-gigabyte su macchine con RAM limitata
\end{itemize}

\subsection{Batch Writing}

La scrittura in batch da 10.000 record bilancia:
\begin{itemize}
    \item \textbf{Overhead I/O}: meno syscall rispetto a scritture singole
    \item \textbf{Memory usage}: batch sufficientemente piccoli da non saturare RAM
\end{itemize}

\subsection{Aggressive Cleanup}

Il pattern di cleanup post-parsing previene memory leak:

\begin{lstlisting}[language=Python]
elem.clear()
while elem.getprevious() is not None:
    del elem.getparent()[0]
\end{lstlisting}

Senza questo cleanup, \texttt{lxml} mantiene riferimenti all'intero documento ancestry.

\chapter{Containerizzazione Docker}

\section{Vantaggi della Containerizzazione}

L'utilizzo di Docker garantisce:

\begin{itemize}
    \item \textbf{Riproducibilità}: ambiente identico su qualsiasi macchina
    \item \textbf{Isolamento}: nessun conflitto con librerie di sistema
    \item \textbf{Portabilità}: deploy immediato su cloud, CI/CD, workstation locali
    \item \textbf{Versionamento}: immagini taggate per rollback facili
\end{itemize}

\begin{tipbox}[Best Practice]
La containerizzazione elimina il classico problema ``works on my machine''. Ogni esecuzione avviene in un ambiente controllato e ripetibile.
\end{tipbox}

\section{Dockerfile}

Il Dockerfile utilizza un'immagine base \texttt{python:3.12-slim} per minimizzare la dimensione:

\begin{lstlisting}[language=docker, caption=Dockerfile]
FROM python:3.12-slim

WORKDIR /app

# Install system dependencies
RUN apt-get update && apt-get install -y \
    build-essential \
    && rm -rf /var/lib/apt/lists/*

# Copy requirements and install
COPY requirements.txt .
RUN pip install --no-cache-dir -r requirements.txt

# Copy source code
COPY . .

# Set python path
ENV PYTHONPATH=/app

# Default command
CMD ["python", "-m", "src.cli"]
\end{lstlisting}

\subsection{Ottimizzazioni Docker}

\begin{description}
    \item[Multi-stage non necessario] Le dipendenze Python non richiedono compilazione pesante
    \item[Cache layer] \texttt{requirements.txt} copiato prima del codice per cache efficiente
    \item[Cleanup apt] Rimozione cache \texttt{/var/lib/apt/lists/*} per immagine più leggera
    \item[--no-cache-dir] Evita cache pip inutile nel container
\end{description}

\section{Docker Compose}

Il file \texttt{docker-compose.yml} semplifica l'esecuzione:

\begin{lstlisting}[language=docker, caption=docker-compose.yml]
services:
  etl:
    build: .
    volumes:
      - ./data:/app/data
      - ./public:/app/public
    environment:
      - PYTHONUNBUFFERED=1
    command: ["tail", "-f", "/dev/null"]
\end{lstlisting}

\subsection{Volume Mounting}

Due volumi persistenti garantiscono che:

\begin{table}[h]
\centering
\begin{tabular}{@{}lp{8cm}@{}}
\toprule
\textbf{Volume} & \textbf{Scopo} \\
\midrule
\texttt{./data:/app/data} & Input XML accessibili dal container \\
\texttt{./public:/app/public} & Output Parquet persistito su host \\
\bottomrule
\end{tabular}
\caption{Volumi Docker}
\end{table}

\section{Workflow Tipico}

\begin{lstlisting}[language=bash, caption=Comandi Docker tipici]
# Build dell'immagine
docker compose build

# Parsing di tutti i file XML
docker compose run --rm etl python -m src.cli parse \
    --input data/ --workers 8

# Query interattiva
docker compose run --rm etl python -m src.cli query \
    --table aiuti --limit 10

# Export CSV
docker compose run --rm etl python -m src.cli export \
    --table aiuti --output public/exports/aiuti.csv
\end{lstlisting}

\section{CI/CD Integration}

La natura containerizzata facilita l'integrazione in pipeline CI/CD:

\begin{figure}[h]
\centering
\begin{tikzpicture}[node distance=1.2cm, scale=0.85, transform shape]
    \node[component, fill=gray!20] (git) {Git Push};
    \node[component, right=of git] (build) {Docker Build};
    \node[component, right=of build] (test) {Run Tests};
    \node[component, right=of test] (deploy) {Deploy};
    
    \draw[arrow] (git) -- (build);
    \draw[arrow] (build) -- (test);
    \draw[arrow] (test) -- (deploy);
\end{tikzpicture}
\caption{Pipeline CI/CD tipica}
\end{figure}

\chapter{CLI e Usabilità}

\section{Framework Click}

L'interfaccia a riga di comando è costruita con \textbf{Click}, un framework Python che offre:

\begin{itemize}
    \item Parsing automatico degli argomenti
    \item Validazione dei tipi
    \item Help auto-generato
    \item Composizione di comandi (gruppi)
\end{itemize}

\section{Comandi Disponibili}

\subsection{parse}

Elabora file XML e li converte in Parquet:

\begin{lstlisting}[language=bash]
docker compose run --rm etl python -m src.cli parse \
    --input data/ \
    --output public/parquet \
    --workers 8
\end{lstlisting}

\begin{table}[h]
\centering
\begin{tabular}{@{}llp{5cm}@{}}
\toprule
\textbf{Opzione} & \textbf{Default} & \textbf{Descrizione} \\
\midrule
\texttt{-i, --input} & (required) & File o directory input \\
\texttt{-o, --output} & \texttt{public/parquet} & Directory output \\
\texttt{-w, --workers} & 4 & Numero processi paralleli \\
\bottomrule
\end{tabular}
\caption{Opzioni comando parse}
\end{table}

\subsection{query}

Esegue query SQL interattive via DuckDB:

\begin{lstlisting}[language=bash]
# Query semplice
docker compose run --rm etl python -m src.cli query \
    --table aiuti --limit 5

# Query SQL custom
docker compose run --rm etl python -m src.cli query \
    --table strumenti \
    --query "SELECT ANNO, SUM(IMPORTO_NOMINALE) as tot 
             FROM read_parquet('public/parquet/strumenti/**/*.parquet') 
             GROUP BY ANNO ORDER BY ANNO"
\end{lstlisting}

\begin{table}[h]
\centering
\begin{tabular}{@{}llp{5cm}@{}}
\toprule
\textbf{Opzione} & \textbf{Default} & \textbf{Descrizione} \\
\midrule
\texttt{-t, --table} & (required) & Tabella target \\
\texttt{-q, --query} & — & Query SQL custom \\
\texttt{-l, --limit} & 10 & Limite risultati \\
\bottomrule
\end{tabular}
\caption{Opzioni comando query}
\end{table}

\subsection{export}

Esporta tabelle in formato CSV o TXT:

\begin{lstlisting}[language=bash]
docker compose run --rm etl python -m src.cli export \
    --table aiuti \
    --format csv \
    --output public/exports/aiuti.csv \
    --delimiter ","
\end{lstlisting}

\subsection{export-aggregated}

Export specializzato con join e aggregazioni pre-calcolate:

\begin{lstlisting}[language=bash]
docker compose run --rm etl python -m src.cli export-aggregated \
    --output public/exports/aggregated.csv
\end{lstlisting}

Questo comando:
\begin{enumerate}
    \item Effettua join tra AIUTI, COMPONENTI e STRUMENTI
    \item Aggrega \texttt{IMPORTO\_NOMINALE} e \texttt{ELEMENTO\_DI\_AIUTO}
    \item Conta componenti e strumenti per aiuto
    \item Produce un file CSV per anno per ottimizzare la memoria
\end{enumerate}

\section{Progress Bar}

Durante l'elaborazione, \texttt{tqdm} fornisce feedback visivo in tempo reale:

\begin{lstlisting}[language=bash, caption=Output con progress bar]
Found 165 XML files to process
Starting processing with 8 workers...
Processing files: 100%|===================| 165/165 [02:15<00:00]
Processing completed in 135.42 seconds
Total processed records: {'aiuti': 2847293, 'componenti': ...}
\end{lstlisting}

\section{Error Handling}

Il sistema gestisce gli errori in modo graceful:

\begin{itemize}
    \item \textbf{File corrotti}: loggati ma non bloccano l'esecuzione
    \item \textbf{failures.txt}: lista dei file falliti salvata per retry
    \item \textbf{Statistiche accurate}: solo file processati correttamente conteggiati
\end{itemize}

\chapter{Inference Service e Classificazione}

\section{Architettura del Microservizio}

Per arricchire il dataset con informazioni relative all'adozione di nuove tecnologie da parte delle imprese, è stato integrato un microservizio di inferenza denominato \textbf{Pack-a-Punch}. Questo componente espone funzionalità di \textit{Natural Language Processing} atte a valutare le descrizioni dei progetti per distinguere le iniziative che impiegano, sviluppano o implementano l'Intelligenza Artificiale, da quelle che non la coinvolgono.

L'architettura dell'inference service si distingue per l'elevato livello di ottimizzazione:
\begin{itemize}
    \item \textbf{Modello Base:} \texttt{dbmdz/bert-base-italian-xxl-cased}. Sfruttando un modello BERT specificamente pre-addestrato sulla lingua italiana, si ottiene un'elevata precisione nell'interpretazione del contesto semantico nazionale.
    \item \textbf{High-Performance Backend:} Il modello allenato (anche tramite knowledge distillation da grandi LLM) è esportato e servito mediante \textbf{ONNX Runtime} con supporto all'accelerazione GPU (CUDA). Ciò abbatte i tempi di inferenza a frazioni di millisecondo per iterazione.
    \item \textbf{Esposizione API:} Sviluppato in \textbf{FastAPI} esponendo l'endpoint REST \texttt{POST /classify}, capace di gestire inferenze in modalità \textit{batch}, riducendo l'overhead di rete tra il layer ETL e il layer di AI.
\end{itemize}

\section{Metodologia di Estrazione e Classificazione}

A causa dell'ingente volume di record (milioni di aiuti), una classificazione "per riga" avrebbe comportato inefficienze strutturali, specialmente constatando l'elevata ripetitività delle descrizioni. La metodologia implementata utilizza un approccio progressivo a \textbf{Multi-livello con Caching}, eseguito in fasi distinte:

\begin{enumerate}
    \item \textbf{Estrazione dei testi univoci:} Sfruttando il motore relazionale \texttt{DuckDB}, vengono interrogate tutte le partizioni Parquet del dataset originale per isolare unicamente le descrizioni distinte nella colonna \texttt{DESCRIZIONE\_PROGETTO}.
    \item \textbf{Fase 1 - Classificazione Binaria:} L'endpoint \texttt{/classify} valuta l'elenco testuale deduplicato con il modello in modalità binaria. L'obiettivo è discernere i progetti strettamente legati all'impiego dell'Intelligenza Artificiale da quelli non pertinenti. I risultati, comprensivi di \textit{confidence score}, popolano una memoria temporanea (\texttt{classification\_cache.parquet}).
    \item \textbf{Fase 2 - Classificazione Multiclasse:} Intervenendo solo sui testi identificati con il tag "AI" al passaggio precedente, il layer ETL li instrada nuovamente verso il servizio di classificazione per individuarne lo specifico ambito applicativo avvalendosi dell'inferenza multiclasse. L'output genera una seconda cache stratificata Parquet (\texttt{multiclass\_cache.parquet}).
    \item \textbf{Aggregazione Finale (Double Left Join):} Infine, le utility contenute nel modulo \texttt{exporter.py} ricongiungono i dataset primari, aggregati parallelamente in chunk annuali con \texttt{ThreadPoolExecutor}, mediante un doppio Join con le due cache pre-costruite. I record passati al modello binario ma che sono risultati non inerenti all'AI confluiranno nell'esportazione CSV presentando i dati multiclasse settati al valore neutro di gestione (NULL / UNKNOWN). 
\end{enumerate}

Questo approccio consolida una metodologia predittiva altamente efficiente: interrogando il modello per la classificazione complessa solamente quando quella binaria accerta il coinvolgimento dell'Intelligenza Artificiale nel progetto, le performance crescono vertiginosamente evitando computazione disperata e irrilevante.

\input{chapters/09-conclusioni}

\end{document}
